\chapter{Conclusiones}

El trabajo desarrollado hasta este punto demuestra la viabilidad de un sistema multiagente para analisis financiero historico basado en una arquitectura modular y trazable. La decision de comenzar con una base determinista ha permitido validar el flujo completo antes de introducir componentes generativos, reduciendo el riesgo de errores dificiles de diagnosticar.

El MVP actual soporta seis intenciones analiticas, gestiona errores basicos de entrada y separa las responsabilidades en familias de agentes. Esta estructura facilita la evolucion del sistema hacia una version hibrida, donde los modelos de lenguaje puedan aportar flexibilidad en la planificacion, generacion de codigo o interpretacion final sin sustituir los mecanismos de validacion y ejecucion controlada.

Como trabajo futuro, se plantea ampliar la evaluacion experimental, comparar el comportamiento de Groq y Gemini en la fase de interpretacion, estudiar la activacion progresiva del LLM en planificacion y analizar bajo que condiciones resulta seguro incorporar generacion de codigo asistida por modelo.
